
\subsection{Incidence projective de jauge du constructeur commun et descente au quotient}

L'image de jauge n'est pas la marginale des liaisons. Son objet complet par niveau est
\begin{equation}\label{eq:pdf2-g-gau-explicito}
 \mathfrak G_{{\rm gau},m}(d_{\rm gau})=
 \bigl(
 \begin{gathered}
  \mathfrak G_{\rm gau}^{\rm int};\\
  \Lambda_m,\widehat X_m,\widehat\mu_m^{\rm ov},
  \{\widehat\mu_{m,g}^{\rm YM},\nu_{m,g},K_{m,g}^{\rm phys}\}_{g>0},\\
  \mathcal G_m^{\rm full},\mathbb P_m^{\rm phys},
  \widehat J_m,\widehat\Gamma_m,
  \{E_{m,c}\}_c,\widehat H_m^{\rm ov},D_m^{\rm YM},\\
  \{\Theta_{\nu,m}\}_{\nu=1}^4,
  \{\mathfrak A_{+,\nu,m},\operatorname{ev}_{t}^{(\nu)},
      \mathcal V_{m,\nu}^{\rm OS}\}_{\nu=1}^4,\\
  \widehat\Omega_m^{(8)},P_{\Omega_m},W_c,
  \operatorname{Carr}_m,\operatorname{Mem}_m,
  \operatorname{Surv}_m,\tau_m,\mathcal P_m^{F^2}
 \end{gathered}
 \bigr).
\end{equation}
Le groupe $\mathcal G_m^{\rm full}$ agit simultanément sur les liens, les repères et les
coordonnées d’incidence $q_c\in\mathbb F_3^6$. La projection physique est sa moyenne
de Haar. Le générateur réduit est
\begin{equation}\label{eq:pdf2-g-gau-generador-estructural}
 D_m^{\rm YM}
 =\left.3\kappa_{\rm av}\sum_c(I-E_{m,c})
  \right|_{\operatorname{Ran}\mathbb P_m^{\rm phys}},
\end{equation}
et $W_c$ appartient à cette sous-algèbre. Les quatre réflexions agissent sur une même mesure, un même noyau et un même terminal.
Le propriétaire exact précédent construit l'action de jauge d'incidence \eqref{eq:amp-ym-accion-gauge-incidencia}, la projection physique de Haar \eqref{eq:amp-ym-proyeccion-fisica-haar} et la compression entrelacée \eqref{eq:amp-ym-compresion-entrelazada}. De cette construction proviennent le noyau complet équivariant $\widehat K_{m,t}^{\rm ov}$, sa descente au quotient physique $K_{m,t}^{\rm phys}$, les quatre opérateurs de réflexion et l'action $U_m(g)$ de $E(4)$ sur le cœur cylindrique lisse. Ce bloc assemble, sans les introduire comme données, le générateur, le vide, le témoin de courbure et les formes OS :
\begin{equation}\label{eq:pdf2-g-gau-entrelazamiento-dato}
 D_{m+1}^{\rm YM}\widehat J_m=\widehat J_mD_m^{\rm YM},
 \quad
 \mathbb P_{m+1}^{\rm phys}\widehat J_m
  =\widehat J_m\mathbb P_m^{\rm phys},
 \quad
 \widehat J_mW_{m,c}=W_{m+1,c}.
\end{equation}
Les première et troisième identités sont la représentation, dans la notation commune, de \eqref{eq:amp-ym-compresion-entrelazada} et \eqref{eq:amp-ym-testigo-fisico-gap} ; le semi-groupe, le vide simple et le saut exact sont déjà obtenus dans \eqref{eq:amp-ym-owner-semigrupo-vacio}--\eqref{eq:amp-ym-owner-brecha-exacta}. Ici, $\nu_{m,g}:=(\pi_m^{\rm phys})_*\widehat\mu_{m,g}^{\rm YM}$ et $K_{m,g}^{\rm phys}(t)=e^{-tD_{m,g}^{\rm YM}}$ dans le quotient physique. La complétude de jauge inclut en outre la désintégration pour chaque collection d'hyperplans euclidiens. Si $A_j\in\mathfrak A_{+,\nu,m}$ et $t_1\le\cdots\le t_n$, la marginale de la même mesure physique satisfait
\begin{equation}\label{eq:pdf2-g-gau-transfer-euclidean-dato}
\begin{aligned}
 &\int\prod_{j=1}^n
   (\tau_{t_je_\nu}A_j)(x)\,d\widehat\mu_{m,g}^{\rm YM}(x)\\
 &\quad=
 \int\prod_{j=1}^nA_j(y_j)\,
  \nu_{m,g}(dy_1)
  \prod_{j=1}^{n-1}K_{m,g}^{\rm phys}
       (t_{j+1}-t_j;y_j,dy_{j+1})\\
 &\quad=
 \left\langle\mathbf1,
  M_{A_1}e^{-(t_2-t_1)D_{m,g}^{\rm YM}}M_{A_2}\cdots
  e^{-(t_n-t_{n-1})D_{m,g}^{\rm YM}}M_{A_n}\mathbf1
 \right\rangle.
\end{aligned}
\end{equation}
La reconstruction précédente prouve l'identité dans les quatre directions au moyen de \eqref{eq:amp-ym-cuatro-os} et construit un unique hamiltonien entrelacé dans \eqref{eq:amp-ym-hamiltoniano-comun}. Par conséquent, la translation OS s'exprime ici par
\begin{equation}\label{eq:pdf2-g-gau-os-generator-dato}
 \mathcal V_{m,\nu}^{\rm OS}T_{m,\nu}^{\rm OS}(s)
 (\mathcal V_{m,\nu}^{\rm OS})^{-1}
 =e^{-sD_{m,g}^{\rm YM}}.
\end{equation}
Ainsi, le générateur des espérances et l'hamiltonien des translations ne sont pas identifiés par leur nom : ce sont deux réalisations déjà construites de la même désintégration, dont la forme limite et la représentation de courbure figurent dans \eqref{eq:amp-ym-forma-limite} et \eqref{eq:amp-ym-f2}.
